% =============================================================================
% Prompt Documentation for severity-eval — EMNLP 2026
% This file documents all evaluation prompts with academic references.
% Use as appendix or supplementary material.
% =============================================================================

% Included from main.tex appendix — section header provided there.
% \section{Evaluation Prompts and Methodology}
% \label{sec:prompts}

We evaluate 22 LLMs across 11 datasets spanning three domains (finance, medical, legal) and two private insurance datasets.
For each dataset, we use two prompt configurations:
\begin{enumerate}
    \item \textbf{Original prompts} (primary results): zero-shot prompts from published evaluation protocols, ensuring comparability with prior work.
    \item \textbf{Standardized prompts} (ablation): uniform HELM-style templates~\cite{liang2023helm} to demonstrate that our findings are robust to prompt choice.
\end{enumerate}

All evaluations use temperature 0.0 (greedy decoding) where supported by the API, and zero-shot prompting (no in-context examples) to avoid confounding few-shot effects with severity-aware evaluation.

% ---------------------------------------------------------------------------
\subsection{Finance Domain}
% ---------------------------------------------------------------------------

\subsubsection{FinanceBench}

\textbf{Source:} Islam et al.~\cite{islam2023financebench}, Table~3 --- Oracle context-first configuration.
The original paper used human evaluation; we add a brief output constraint (``Provide only the answer'') necessary for automated scoring, following standard practice in benchmark evaluation~\cite{liang2023helm}.

\begin{quote}
\texttt{Answer this question: \{question\}} \\
\texttt{Provide only the answer, no explanation.} \\
\texttt{Context:} \\
\texttt{[START OF FILING] \{evidence\} [END OF FILING]} \\
\texttt{Answer:}
\end{quote}

Evidence is extracted from the \texttt{evidence} column of the dataset, containing SEC filing excerpts (10-K, 10-Q, 8-K) linked to each question.
Parameters: temperature 0.01, max tokens 2048 in the original paper; we use temperature 0.0, max tokens 512.

\subsubsection{FinQA}

\textbf{Source:} Chen et al.~\cite{chen2022program} (Program-of-Thoughts), confirmed by the PIXIU/FinBen benchmark~\cite{xie2024finben} (Appendix~C, Table~7).

\begin{quote}
\texttt{Read the following text and table, and then answer a question:} \\
\texttt{\{evidence\}} \\
\texttt{Question: \{question\}} \\
\texttt{Answer:}
\end{quote}

Context is assembled from the raw JSON: \texttt{pre\_text} paragraphs (joined), followed by the table (pipe-delimited rows), followed by \texttt{post\_text} paragraphs.
This matches the CoT format from~\cite{chen2022program} without the Python code generation component, as we evaluate direct answer generation.

\subsubsection{TAT-QA}

\textbf{Source:} PIXIU/FinBen benchmark~\cite{xie2024finben}, Appendix~C, Table~7; also confirmed by TAT-LLM~\cite{zhu2024tatllm}, Appendix~D.

\begin{quote}
\texttt{Read the following text and table, and then answer a question:} \\
\texttt{\{evidence\}} \\
\texttt{Question: \{question\}} \\
\texttt{Answer:}
\end{quote}

Context is assembled from the passage-level JSON: table rows (pipe-delimited) followed by paragraph texts.
The FinBen benchmark uses the same structure with the prefix ``Please answer the given financial question based on the context''~\cite{xie2024finben}; we adopt the FinQA-consistent format for uniformity across financial datasets.

% ---------------------------------------------------------------------------
\subsection{Medical Domain}
% ---------------------------------------------------------------------------

\subsubsection{HEAD-QA}

\textbf{Source:} MedHELM~\cite{bedi2025medhelm}\footnote{\url{https://github.com/stanford-crfm/helm/blob/main/src/helm/benchmark/run_specs/medhelm_run_specs.py}} --- zero-shot biomedical MCQ format.

\begin{quote}
\texttt{You are a highly knowledgeable AI assistant specializing in biomedical sciences.
Select the correct answer by outputting only the letter corresponding to your choice.} \\[6pt]
\texttt{Question: \{question\}} \\
\texttt{A. \{option\_A\}} \\
\texttt{...} \\
\texttt{Answer:}
\end{quote}

The biomedical instruction prefix follows MedHELM's domain-specific adapter configuration.
The original HEAD-QA paper~\cite{vilares2019headqa} evaluated IR+neural systems without LLM prompts; the lm-evaluation-harness uses a simpler format without the biomedical prefix.

% ---------------------------------------------------------------------------
\subsection{Legal Domain}
% ---------------------------------------------------------------------------

All legal prompts follow LegalBench~\cite{guha2023legalbench}, the standard benchmark for LLM evaluation on legal tasks.

\subsubsection{CUAD}

\textbf{Source:} LegalBench~\cite{guha2023legalbench} --- zero-shot jargon format.
The original CUAD paper~\cite{hendrycks2021cuad} defines an extractive QA task (SQuAD~2.0 format) where questions are the 41 clause category descriptions from \texttt{category\_descriptions.csv}.

\begin{quote}
\texttt{\{evidence\}} \\[6pt]
\texttt{Question: \{question\}} \\
\texttt{Provide only the relevant text from the clause, or `N/A' if not found.} \\
\texttt{Answer:}
\end{quote}

Evidence is the contract paragraph context from the SQuAD-format data, downloaded directly from the GitHub repository\footnote{\url{https://github.com/TheAtticusProject/cuad}} since the HuggingFace dataset script is no longer supported.

\subsubsection{MAUD}

\textbf{Source:} LegalBench~\cite{guha2023legalbench} --- MCQ classification format over merger agreement text.
MAUD~\cite{wang2023maud} defines 92 deal point questions with 2--5 answer categories per question.

\begin{quote}
\texttt{Instruction: Read the segment of a merger agreement and answer the
multiple-choice question by choosing the option that best characterizes the agreement.} \\
\texttt{Question: \{question\}} \\
\texttt{A. \{option\_A\}} \\
\texttt{B. \{option\_B\}} \\
\texttt{...} \\[6pt]
\texttt{Merger Agreement: \{evidence\}} \\
\texttt{Answer:}
\end{quote}

Answer options are pre-computed per question type from the full dataset.
The LegalBench evaluation includes both few-shot (\texttt{base\_prompt.txt}) and zero-shot (\texttt{base\_prompt\_wo\_example.txt}) variants; we use the zero-shot variant.

\subsubsection{ContractNLI}

\textbf{Source:} Adapted from LegalBench~\cite{guha2023legalbench} ContractNLI tasks.
The original paper~\cite{koreeda2021contractnli} defines a document-level NLI task with 17 fixed hypotheses over NDA contracts.
LegalBench binarizes the task (Entailment vs.\ Not Entailment); we preserve the original three-class formulation (Entailment, Contradiction, Neutral).

\begin{quote}
\texttt{Given the following NDA excerpt, determine the relationship with the hypothesis.
Answer with exactly one word: Entailment, Contradiction, or Neutral.} \\[6pt]
\texttt{Document excerpt:} \\
\texttt{\{evidence\}} \\[6pt]
\texttt{Hypothesis: \{question\}} \\
\texttt{Answer:}
\end{quote}

% ---------------------------------------------------------------------------
\subsection{Insurance Domain (Private Datasets)}
% ---------------------------------------------------------------------------

\subsubsection{Quebec Automobile Insurance RAG}

\textbf{Source:} Beauchemin et al.~\cite{beauchemin2024quebec} --- exact prompt from the RAG evaluation pipeline.

\begin{quote}
\texttt{Répondez à la question suivante sur l'assurance automobile au Québec.
Fournissez uniquement la réponse, sans explication.} \\[6pt]
\texttt{Question : \{question\}} \\
\texttt{Réponse :}
\end{quote}

The original paper uses a more detailed expert system prompt with retrieved context; we use the zero-shot variant for consistency.
Evaluation follows the manual criterion-based scoring protocol from~\cite{beauchemin2024quebec} (scale: $-1$ for false information, 0 for no answer, 1 for incomplete, 2 for complete).

\subsubsection{Legal Text Simplification (JUDGEBERT)}

\textbf{Source:} JUDGEBERT~\cite{judgebert2025} --- GPT-4 simplification prompt (French).

\begin{quote}
\texttt{Simplifiez le texte juridique suivant en langage clair,
en préservant le sens légal.} \\[6pt]
\texttt{Texte : \{question\}} \\
\texttt{Simplification :}
\end{quote}

Evaluation of simplification quality uses the Legal Meaning Preservation (LMP) score from~\cite{judgebert2025}, a 10-point Likert scale assessed by law students with brackets: 7--10 (Accurate), 2--6 (Seems Imprecise), 1 (Off-Track).

% ---------------------------------------------------------------------------
\subsection{Standardized Prompts (Ablation)}
\label{sec:standard-prompts}
% ---------------------------------------------------------------------------

For the ablation study, all datasets use a uniform HELM-style template~\cite{liang2023helm}:

\begin{quote}
\texttt{Question: \{question\}} \\
\texttt{\{context if available\}} \\
\texttt{\{options if MCQ\}} \\
\texttt{Answer:}
\end{quote}

This minimal template removes all dataset-specific instructions, domain prefixes, and formatting conventions.
If the Kendall $\tau$ between model rankings under ``original'' and ``standard'' prompts is high ($\tau > 0.8$), this confirms that our severity-weighted findings are robust to prompt engineering choices.

% ---------------------------------------------------------------------------
\subsection{Severity Annotation}
\label{sec:severity-annotation}
% ---------------------------------------------------------------------------

Each dataset instance is annotated with a severity level (critical, major, minor, negligible) using domain-specific heuristics:

\paragraph{Finance (FinanceBench, FinQA, TAT-QA).}
Two-axis classification following SEC Staff Accounting Bulletin No.~99 (SAB~99) materiality framework:
(1)~answer type (dollar-large, dollar-small, percentage, ratio, yes/no, text) $\times$
(2)~metric type (core financials, operational, ratio, descriptive).
Core financial statement items (revenue, net income, total assets) in absolute dollars are classified as \emph{critical}; operational metrics as \emph{major}; ratios and percentages as \emph{minor}; descriptive text as \emph{negligible}.

\paragraph{Medical (HEAD-QA).}
Based on clinical impact of erring on the instance: pharmacology
questions are \emph{critical}; medicine and nursing are \emph{major};
psychology is \emph{minor}; biology and chemistry are
\emph{negligible}.

\paragraph{Legal (CUAD, MAUD, ContractNLI).}
Based on clause type and legal exposure: indemnification, liability limitation, and termination clauses are \emph{critical}; IP ownership, non-compete, and deal protection are \emph{major}; governing law and closing conditions are \emph{minor}; definitions and boilerplate are \emph{negligible}.

% ---------------------------------------------------------------------------
\subsection{Scoring Methodology}
% ---------------------------------------------------------------------------

Predictions are scored using an automated cascade:
\begin{enumerate}
    \item \textbf{MCQ matching}: if options are provided, compare by option letter or exact text match.
    \item \textbf{Yes/No extraction}: word-boundary match at response start.
    \item \textbf{Numeric comparison}: extract dollar amounts (priority), percentages, or first non-year number; compare with 5\% relative tolerance.
    \item \textbf{Exact text match}: normalized (lowercase, strip punctuation, collapse whitespace).
    \item \textbf{Fuzzy containment}: reference substring in prediction (min 4 characters) or all reference words present (min 3 words).
\end{enumerate}

Number extraction prioritizes dollar amounts (\$X, \$X million/billion) over standalone numbers, and skips year-like values (1900--2099) to avoid false matches in financial documents.

% References are handled by the main bibliography (references.bib).
